\section*{Chapter \ref{ch:iv1:machine-learning} --- Machine Learning}
\markboth{\protect\omTocMark{}{Answers: Machine Learning}}{\protect\omTocMark{}{Answers: Machine Learning}}

\begin{solution}{iq:iv1:machine-learning:1}
Accuracy $0.93$, precision $0.40$, recall $0.80$, F1 $0.53$. A model that never flags anything has accuracy $0.95$, which is higher.

The counts are 40 true positives, 60 false positives, 10 false negatives, and so 890 true negatives. Accuracy is $930/1000$. Precision\index{precision and recall} is $40/100$, the share of flags that are right. Recall is $40/50$, the share of real events that are caught. F1\index{F1 score} is their harmonic mean, $2(0.4)(0.8)/1.2 = 0.533$.

The empty model is right on the 950 days without the event, so it beats the real model on accuracy while being useless. When classes are rare, accuracy measures the base rate.

\iqwrong{Reporting accuracy of $0.93$ as a strong result. Or adding 40 and 10 to get the flagged days, which confuses the two denominators.}

\iqfollow{Write F1 with no reference to true negatives.}{$\mathrm{TP}/(\mathrm{TP} + \frac12(\mathrm{FP} + \mathrm{FN})) = 40/75$. The 890 true negatives have no effect, which is the point of the measure.}

\iqwhy{This is the fastest check that a candidate reads the right denominators, and the example with a useless model shows whether accuracy will be questioned.}
\end{solution}

\begin{solution}{iq:iv1:machine-learning:2}
About 59\% accuracy, though every model is worth 50\%. A test set is used once, after all selection is done.

Each model's score is roughly normal, with mean 50\% and standard deviation 5\%. The largest of 20 independent normals lies about $1.87$ standard deviations above the mean, so the expected best is about $59.3\%$, and an exact binomial calculation and a simulation give the same value.

The right procedure: choose on a separate validation set, or by cross-validation\index{cross-validation}, and then score the single chosen model once on the untouched test set. That score is unbiased for the chosen model, and here it would come back near 50\%. If several models must be reported, adjust for the number tried.

\iqwrong{Reporting the best model's test score, or using the test set to decide what to try next. Each peek turns the test set into a validation set.}

\iqfollow{How long a test set keeps the best of 20 null models below 52\% on average?}{About 2200 days: $1.87 \times 50/\sqrt{n} = 2$ gives $n \approx 2180$.}

\iqwhy{The whole subject of overfitting research can be put in this one example. Quant interviewers listen for whether the candidate counts trials without being asked.}
\end{solution}

\begin{solution}{iq:iv1:machine-learning:3}
Ridge, $k$-nearest neighbours and principal component regression change. Ordinary least squares, the tree and the forest do not, apart from any random seeds.

Least squares rescales the coefficient by $1/100$, and predictions are unchanged. Ridge\index{ridge regression} penalises the size of coefficients, so a smaller coefficient on the rescaled feature pays less, and the fit shifts. Nearest neighbours uses distances, so the feature in cents dominates them. Principal components chase variance, so the same feature takes over the first component.

A tree splits on order alone, and multiplying by 100 preserves order; a forest is a collection of trees. Runs on simulated data confirm each case.

\iqwrong{Saying that ridge is unaffected because least squares is. Or that all models should be fine after fitting. The rule is: anything that uses a distance or a penalty on size needs a common scale.}

\iqfollow{What should you do with such models, and when?}{Standardise each feature, using the training data only. Using all the data leaks information from the test period.}

\iqwhy{A basic preparation question that catches people who treat all models alike, and it is a real source of bugs.}
\end{solution}

\begin{solution}{iq:iv1:machine-learning:4}
The constant $0.1$ earns $0.325$, the constant $0.5$ earns $0.693$, and the single row costs $\ln 1000 = 6.91$.

Log loss\index{log loss} is the mean of $-\ln(\text{predicted probability of what happened})$. For the constant $0.1$: $-0.1\ln 0.1 - 0.9\ln 0.9 = 0.325$. For $0.5$: $-\ln 0.5 = 0.693$, for every row, whatever the base rate.

A prediction of $0.999$ on a label of 0 gives probability $0.001$ to what occurred, and it costs $6.91$. That is more than twenty times the base loss, and one such row in 100 adds $0.069$, a fifth of the base loss. So the metric punishes overconfidence severely. A model that reports raw scores near 0 and 1 needs clipping or calibration.

\iqwrong{Comparing a model's log loss to $0.693$ on an unbalanced problem. The fair baseline is the constant at the base rate, here $0.325$.}

\iqfollow{Which constant minimises log loss?}{The base rate. Setting the derivative of $-b\ln p - (1 - b)\ln(1 - p)$ to zero gives $p = b$, which is why log loss rewards calibration.}

\iqwhy{Reading a loss needs a baseline. A candidate who gives the base-rate figure at once has used this metric on real data.}
\end{solution}

\begin{solution}{iq:iv1:machine-learning:5}
The odds double, since $e^{0.7} = 2.01$. A probability of $0.2$ becomes $0.33$, and $0.5$ becomes $0.67$.

The model is linear in the log-odds\index{log-odds}, so a unit rise adds $0.7$ to the log-odds and multiplies the odds by $e^{0.7}$.

From $0.2$, the odds are $0.25$, which double to $0.5$, and so $p = 0.5/1.5 = 1/3$. From $0.5$, the odds are 1, which double to 2, and so $p = 2/3$. The same coefficient adds $0.13$ in one case and $0.17$ in the other. The effect on the probability depends on where you start.

\iqwrong{Saying that the probability rises by $0.7$, or doubles. A probability cannot do that from $0.6$, whereas the odds can.}

\iqfollow{Give a quick rule for the effect near $p = \frac12$.}{The slope of the logistic curve there is $\frac14$, so the change is about the coefficient over 4. Here that is $0.175$, close to the true $0.17$.}

\iqwhy{Coefficients are read wrongly more often than they are computed wrongly. Separating odds from probabilities is what the question tests.}
\end{solution}

\begin{solution}{iq:iv1:machine-learning:6}
Training sets have 800 rows. A row is validated once and helps train four models. The search needs 5000 fits.

Each fold holds out one fifth, 200 rows, and trains on the other 800. Over five folds, a row is in the held-out part exactly once and in the training part the other four times.

The grid has $10^3 = 1000$ points, and each needs 5 fits, for 5000 in all\index{hyperparameter search}. Once the point is chosen, one more fit on all 1000 rows gives the final model.

The cross-validated score of the chosen point is biased upward, since it is the best of 1000. A separate test set, or an outer loop, is needed for an honest figure.

\iqwrong{Reporting the score of the best grid point as the expected performance, or multiplying 3 by 10 for the size of the grid. It is 10 to the power 3.}

\iqfollow{What does a nested search with 5 outer folds cost?}{$25\,000$ fits, five times as many. That is what an unbiased estimate of the whole procedure costs.}

\iqwhy{Simple counts reveal who has actually run such a search, and who has thought about what its best score means.}
\end{solution}

\begin{solution}{iq:iv1:machine-learning:7}
About 15\% of flags are real, 48\% at the lower false positive rate, and the ROC area does not change with prevalence.

Per 1000 trades: 10 are bad, of which 9 are caught, and 990 are good, of which $49.5$ are flagged. So precision is $9/58.5 = 0.154$. With a false positive rate of 1\%, the flags are $9 + 9.9$, which gives $0.476$. A small improvement in the false positive rate has a large effect, because the good trades outnumber the bad.

The ROC curve plots the true positive rate against the false positive rate\index{ROC curve}. Each is computed within a single class, so the mix of classes does not enter. A run at both prevalences gives the same area.

\iqwrong{Saying that a detector with 90\% and 95\% rates is right most of the time when it flags. Or reporting ROC area alone for a rare-event problem, where it hides the false alarm load.}

\iqfollow{What is the share of real flags at 10\% prevalence?}{$0.09/(0.09 + 0.045) = 0.667$, with the same detector. Precision does depend on prevalence, so a figure measured on a balanced test set does not carry over to production.}

\iqwhy{It separates metrics that depend on prevalence from those that do not, and it comes up whenever a model moves to a stream with a different mix.}
\end{solution}

\begin{solution}{iq:iv1:machine-learning:8}
There are 3457 parameters, about $2.8$ per row, so with low-signal data the model will memorise unless heavily constrained.

Each layer has inputs times units for the weights, plus one bias per unit. Layer one: $(20 + 1) \times 64 = 1344$. Layer two: $(64 + 1) \times 32 = 2080$. Output: $(32 + 1) \times 1 = 33$. The sum is 3457, and the count from a deep learning library agrees.

With 1250 rows, the model has more parameters than rows, and returns are mostly noise, so it can fit the training set with no real signal. Weight decay, dropout, early stopping and a smaller architecture are the normal responses. An honest walk-forward validation decides. A very small network, or a linear model, is the baseline to beat.

\iqwrong{Forgetting the biases, which gives 3360. Or concluding that the network must fail: overparameterised networks can generalise, but on noisy returns the size of the model is a warning that needs testing.}

\iqfollow{What would you try first with so little data?}{A regularised linear model, as the baseline. If a network cannot beat it out of sample, the extra parameters are not paying.}

\iqwhy{Counting parameters is mechanical, and what the interviewer listens for is the comparison with the number of rows.}
\end{solution}

\begin{solution}{iq:iv1:machine-learning:9}
Model A overfits, with high variance. Model B underfits, with high bias. More data helps A and not B.

A has a gap of 29 points between training and validation. It is memorising. More data narrows the gap\index{learning curves}, as do stronger regularisation, fewer features and a simpler model.

B has almost no gap, and both are far below the target of 85\%. It lacks capacity or information. Adding rows moves both curves very little. It needs a more flexible model, or better features. If the target is out of reach for any model, the noise in the labels may be the real limit, so check this first.

A run shows both. A 1-nearest-neighbour classifier on a noisy wave has a wide gap that narrows with more rows. A linear classifier on the same wave stays low on both sets and does not improve.

\iqwrong{Saying that more data always helps. It does nothing for a model that is already too simple, as B's curves have met already.}

\iqfollow{What would a validation score higher than the training score suggest?}{Regularisation active only in training, such as dropout, or an easier validation period. If neither applies, suspect a split that leaks.}

\iqwhy{Diagnosis decides where time goes. A candidate who can say which fix is wasted saves weeks of work.}
\end{solution}

\begin{solution}{iq:iv1:machine-learning:10}
Bias is $0.1(k + 1)$ and variance $1/k$, so the best $k$ is 3, with squared error $0.493$. Doubling the density raises the best $k$ to 6.

The mean of the $k$ neighbours has expected value $0.2(1 + \dots + k)/k = 0.1(k + 1)$ away from the true value\index{bias and variance}. Its variance is $1/k$. So the squared error is
\[
\frac{1}{k} + 0.01(k + 1)^2.
\]
The values at $k = 1, 2, 3, 4$ are $1.04$, $0.59$, $0.493$ and $0.5$, so the best is $k = 3$. A simulation of the three-neighbour estimate gives the same $0.493$.

With spacing $0.1$, the error is $1/k + 0.0025(k + 1)^2$, with minimum $0.289$ at $k = 6$. More data allows a larger neighbourhood, which lowers the variance while the distances stay small. Both bias and variance fall.

Each new observation still has noise of variance 1, which no choice of $k$ removes.

\iqwrong{Adding the bias to the variance with no square. The terms that add are the variance and the squared bias.}

\iqfollow{If the neighbours lie evenly on both sides, what is the best $k$?}{As large as the data allow, in this model. A straight line gives zero bias under symmetric averaging, so only the variance remains. Curvature is what would stop it.}

\iqwhy{An explicit bias-variance curve is more convincing than the usual picture, and it shows how the best complexity moves with the data.}
\end{solution}

\begin{solution}{iq:iv1:machine-learning:11}
Training error $0.8$, new-noise error $1.2$. The gap of $0.4$ equals $2p\sigma^2/n$.

Fitting projects the noise on the regressors' column space. The fit absorbs $p$ of the $n$ dimensions of noise, so the residual sum of squares has expectation $\sigma^2(n - p)$\index{optimism}, and the mean squared error is $1 - 20/100 = 0.8$.

The fitted values carry noise of total variance $p\sigma^2$. Against fresh noise at the same points, the expected error is the new noise plus that: $1 + 20/100 = 1.2$. A simulation returns $0.80$ and $1.20$.

The fit therefore looks better than the noise level, and predicts worse by the same amount. The gap is linear in the number of parameters. This is what Mallows's $C_p$ and the Akaike criterion add back to the training error.

\iqwrong{Saying that a correct model has no overfitting. Being correct removes the bias, not the variance from estimating 20 numbers.}

\iqfollow{What do ten more useless parameters add to the gap?}{$2 \times 10/100 = 0.2$. A criterion accepts a new feature only if it lowers the training error by more than this charge.}

\iqwhy{This is the simplest exact result behind model selection. It shows that the training error misleads even when the model is right.}
\end{solution}

\begin{solution}{iq:iv1:machine-learning:12}
The Gini impurity falls by $0.163$, and the entropy by $0.256$ bits.

Left child: 30 and 10, right child: 10 and 50, with weights $0.4$ and $0.6$.

Gini\index{Gini impurity} is $2p(1 - p)$. Parent: $2(0.4)(0.6) = 0.48$. Left: $2(0.75)(0.25) = 0.375$. Right: $2(\frac16)(\frac56) = 0.278$. Weighted: $0.4(0.375) + 0.6(0.278) = 0.317$. So the decrease is $0.163$.

Entropy\index{entropy} is $-p\log_2 p - (1 - p)\log_2(1 - p)$. Parent: $0.971$. Left: $0.811$. Right: $0.650$. Weighted: $0.715$. The decrease is $0.256$.

The two criteria are on different scales, so the numbers are not comparable, and in practice they rank most splits the same way.

\iqwrong{Averaging the two children with equal weights, which gives a different number. The weights are the shares of rows.}

\iqfollow{What is the meaning of the Gini impurity?}{The chance that two rows drawn at random with replacement from the node have different labels, or the error of a rule that guesses at random in the node's proportions.}

\iqwhy{A simple check that someone can compute what a library computes, with the weighting as the common mistake.}
\end{solution}

\begin{solution}{iq:iv1:machine-learning:13}
Misclassification error gives $0.25$ for both. Gini gives $0.375$ for A and $0.333$ for B, so B is preferred, since it creates a pure node. The cause is that Gini is strictly concave.

Under the majority rule: in A, each child is wrong on 100, for $200/800$. In B, the first child is wrong on 200 and the second on 0, also $200/800$.

Gini for A: $2(\frac34)(\frac14) = 0.375$ for each child. For B: the first child has $\frac34$ of the rows, with $2(\frac13)(\frac23) = 0.444$, and the pure child has 0, so $0.75 \times 0.444 = 0.333$.

The misclassification rate is piecewise linear in $p$\index{split criterion}. Unless the majority changes across the children, mixing two children gives the same weighted error, so the split shows no gain. A strictly concave function always falls under a split that separates the classes at all, and it rewards purity. The pure node is valuable since later splits only have to work on the rest.

\iqwrong{Saying that A is better because it is symmetric, or concluding that the criteria always disagree. They usually agree, but ties and near-ties, where they can differ, are frequent early in a tree. Children $(210, 400)$ and $(190, 0)$, for instance, have a worse error than A, $0.2625$, and a better Gini, $0.344$.}

\iqfollow{When does a split give zero gain, for every concave criterion?}{When both children have the parent's proportions. In that case nothing has been learned about the label.}

\iqwhy{It explains why a tree does not simply optimise the metric used in the end, and that is a useful general principle.}
\end{solution}

\begin{solution}{iq:iv1:machine-learning:14}
The variance is $\rho\sigma^2 + (1 - \rho)\sigma^2/B$, which is $0.307$ at $B = 100$, and $0.3$ in the limit. The trade is good, with $0.131$ at the same $B$.

Expand the variance of the mean\index{bagging}:
\[
\frac{1}{B^2}\left(B\sigma^2 + B(B - 1)\rho\sigma^2\right) = \rho\sigma^2 + \frac{(1 - \rho)\sigma^2}{B}.
\]
At the stated values: $0.3 + 0.7/100 = 0.307$. Check the limits: $\rho = 1$ gives $\sigma^2$, as one model does, and $\rho = 0$ gives $\sigma^2/B$. A simulation of correlated models returns $0.307$.

So more models stop helping at the shared part $\rho\sigma^2$. To go lower, the models must be made less alike, and that is the purpose of feature subsampling in a random forest\index{random forest}. With $\rho = 0.1$ and $\sigma^2 = 1.2$: $0.12 + 1.08/100 = 0.131$. The individual models are worse, while the average is more than twice as good.

\iqwrong{Assuming that the variance goes to zero with enough trees. That holds only for independent models, which trees built on the same data are not.}

\iqfollow{At $\rho = 0.3$, how many models get within 1\% of the limit?}{$0.7/B \le 0.003$ gives about 234. Past this, adding more is waste.}

\iqwhy{The formula explains why forests work and why ensembles of similar alphas disappoint. Quants meet the same formula when they combine correlated signals.}
\end{solution}

\begin{solution}{iq:iv1:machine-learning:15}
The fitted share after $m$ trees is $1 - 0.9^m$, which passes 90\% at $m = 22$ and 99\% at $m = 44$. Shrinking averages over the noise.

For $L = \frac12(y - F)^2$, the derivative with respect to $F$ is $-(y - F)$. So the negative gradient is the residual, and fitting a tree to residuals is a gradient step in function space\index{gradient boosting}.

If each tree fits the residual exactly and is added with weight $\nu$, the residual is multiplied by $1 - \nu = 0.9$ every round. After $m$ rounds, $0.9^m$ remains. Since $0.9^{22} = 0.098$ and $0.9^{44} = 0.0097$, the two thresholds are crossed at 22 and 44.

A real tree fits a noisy estimate of the residual. With $\nu = 1$, the whole of the noise in the last tree stays in the fit. With a small $\nu$, the fit is an average over many trees, each with different noise. The price is more trees.

\iqwrong{Thinking that shrinkage changes the final target. It only changes how fast it is reached and how much noise comes along.}

\iqfollow{If each tree's estimate of the residual has noise of variance 1, independently, what noise variance does the fit settle at?}{$\nu/(2 - \nu)$, which is $0.053$ at $\nu = 0.1$ and 1 at $\nu = 1$. A simulation agrees.}

\iqwhy{A candidate should know what boosting does at the level of one row. The formula explains the usual advice about small learning rates with more trees.}
\end{solution}

\begin{solution}{iq:iv1:machine-learning:16}
The forest's test error levels off and stays there, a little above the 20\% noise. The boosting training error goes to zero, while its test error falls and then rises.

A forest averages trees. Adding a tree changes only the estimate of an average, whose variance falls with $B$, so the error approaches a limit. It cannot overfit by having too many trees\index{random forest}, although each deep tree has fitted the noise, and the forest's training error is close to zero too.

Boosting is sequential\index{gradient boosting}. Each tree aims at what is still unexplained, and once the pattern is exhausted, this is the flipped labels. The model has unbounded capacity in the number of rounds, and it goes on fitting them.

A run on 400 rows, with a checkerboard pattern and 20\% flips, shows it. Boosting with learning rate $0.5$ reaches a training error below 1\%, and its test error ends above 25\%, more than two points above its minimum. For the forest, going from 100 to 1000 trees moves test error by less than a point. So in boosting the number of trees is a hyperparameter to tune by early stopping, and in a forest it is a question of computation.

\iqwrong{Saying that both overfit with more trees, or that a forest never overfits. It does not overfit through the number of trees, but it can through deep trees on noisy data.}

\iqfollow{How do you choose the number of rounds on financial data?}{By early stopping on a later block of time, separated from the training data by a gap, and not on a random subset of the rows.}

\iqwhy{It is the central contrast between the two families. The interviewer wants the reason, not only the fact.}
\end{solution}

\begin{solution}{iq:iv1:machine-learning:17}
$\alpha = 0.424$. Wrong rows are multiplied by $1.528$ and right rows by $0.655$. After normalising, the wrong rows hold a total weight of exactly $\frac12$.

$\alpha = \frac12\ln(0.7/0.3) = 0.424$\index{AdaBoost}. The factors are $e^{\alpha} = \sqrt{7/3}$ and $e^{-\alpha} = \sqrt{3/7}$.

The wrong rows held $0.3$, so they now hold $0.3\sqrt{7/3} = \sqrt{0.21}$. The right rows held $0.7$, now $0.7\sqrt{3/7} = \sqrt{0.21}$. These are equal, so each share is $\frac12$. The normaliser is $2\sqrt{0.21} = 0.917$.

This means that the last learner now has a weighted error of exactly $\frac12$ on the new weights: it is worth nothing. So the next learner must find something new. It cannot score by repeating the last one. That holds whatever $\varepsilon$ is, which is the reason for this choice of $\alpha$.

\iqwrong{Saying that wrong rows get ``more weight'' with no number. The exact $\frac12$ is the content of the algorithm.}

\iqfollow{What are the votes for $\varepsilon = 0.5$ and $\varepsilon = 0.6$?}{Zero and $\frac12\ln(2/3) = -0.20$. A coin flip gets no vote, and a learner worse than a coin gets a negative vote, so that its opposite is used.}

\iqwhy{It tests whether the reweighting is known as a calculation, and it leads to the bound on the training error.}
\end{solution}

\begin{solution}{iq:iv1:machine-learning:18}
$\lambda = \sigma^2/\tau^2$. The effective number of parameters is $\sum d_i^2/(d_i^2 + \lambda) = 1.5$.

The log posterior is, up to a constant,
\[
-\frac{1}{2\sigma^2}\|y - X\beta\|^2 - \frac{1}{2\tau^2}\|\beta\|^2.
\]
Multiplying by $2\sigma^2$ gives the ridge objective\index{ridge regression}, with $\lambda = \sigma^2/\tau^2$. The posterior is normal, so its mode is its mean, $(X^\top X + \lambda I)^{-1}X^\top y$. A tighter prior or noisier data means more shrinkage.

The fit is $\hat{y} = Hy$ with $H = X(X^\top X + \lambda I)^{-1}X^\top$. The natural definition of effective parameters\index{effective degrees of freedom} is the trace of $H$, since for least squares this equals $p$. In terms of the singular values $d_i$ of $X$, it is $\sum d_i^2/(d_i^2 + \lambda)$.

With values $100, 10, 1$ and $\lambda = 10$: $100/110 + 10/20 + 1/11 = 0.909 + 0.5 + 0.091 = 1.5$. A matrix with these singular values gives the same trace. At $\lambda \to 0$ the sum is 3, the full count, and at $\lambda \to \infty$ it is 0. The strong direction keeps its freedom, and the weak one is nearly removed.

\iqwrong{Treating $\lambda$ as a tuning knob with no units. It depends on the scale of the features and on the noise, which the Bayesian form makes obvious.}

\iqfollow{What is $\lambda$ for noise of 1 and a prior standard deviation of $0.1$?}{$\lambda = 1/0.01 = 100$. For returns, where signals are small, large values of $\lambda$ are normal and not a sign of a mistake.}

\iqwhy{The two views of ridge, penalty and prior, are often asked together. Effective parameters are what let a candidate compare a ridge fit with an unpenalised one.}
\end{solution}

\begin{solution}{iq:iv1:machine-learning:19}
With no penalty, $\beta$ runs off to infinity. With $\lambda = 0.1$, $\beta \approx 2.13$, and the probability at $x = 1$ is about $0.89$.

The loss is $2\ln(1 + e^{-\beta})$, which falls forever as $\beta$ rises, and approaches 0 with no minimiser. Separable data\index{separable data} has no finite maximum likelihood estimate. In practice the coefficient grows until the solver stops, and the software may warn or may not.

With the penalty, set the derivative to zero:
\[
-\frac{2}{1 + e^{\beta}} + 0.1\beta = 0.
\]
That equation has a single root at $\beta = 2.13$, found by bisection and matched by a library fit with the same penalty. The probability is $1/(1 + e^{-2.13}) = 0.894$.

So the penalty has a direct meaning here: it decides how confident the model may be with no contradicting data. The result is a finite and sensible probability from a dataset of two rows.

\iqwrong{Reporting the unpenalised fit as $\beta$ large but finite, or as a probability of 1. Also saying that separable data is a good sign. With many features and few rows, it is often just a sign of overfitting.}

\iqfollow{What happens at $\lambda = 0.01$?}{$\beta \approx 3.91$, with probability $0.98$. Weaker penalties allow more confidence, slowly: a tenfold fall in $\lambda$ does not even double $\beta$.}

\iqwhy{This is a classic numerical surprise. The candidate should know why a logistic fit fails to converge, and what regularisation does about it.}
\end{solution}

\begin{solution}{iq:iv1:machine-learning:20}
The true posterior is $0.23$ and the classifier says $0.73$, so the decision flips.

Naive Bayes\index{naive Bayes} multiplies the likelihood ratios, assuming independence given the class. With prior $\frac12$, the posterior odds are the product of the ratios.

The true odds use A once: $3 \times 0.1 = 0.3$, so the posterior is $0.3/1.3 = 0.23$. Below $\frac12$, the decision is class 0.

The classifier counts A three times: $3^3 \times 0.1 = 2.7$, so the posterior is $2.7/3.7 = 0.73$. Above $\frac12$, the decision is class 1. A true evidence ratio of 3 has been treated as 27.

With A alone, the problem is hidden: the true posterior is $0.75$, and the classifier's is $0.964$. These fall on the same side of $\frac12$, so the error is only in confidence. A second feature is what turns it into a wrong decision.

\iqwrong{Saying that naive Bayes is immune to correlation because it often ranks well. That is often true of ranking, but correlated features cause overconfidence that can change decisions.}

\iqfollow{Would logistic regression suffer in the same way?}{No. It shares the coefficient among the copies, so that their sum stays the same, and it returns the same fit. A run with three copies confirms this.}

\iqwhy{It makes concrete what the independence assumption costs. In finance, where features are highly correlated, this is an everyday issue.}
\end{solution}

\begin{solution}{iq:iv1:machine-learning:21}
Both gradients are $0.5$, the new weights are $-0.05$ and $1.95$, and the loss falls from $0.5$ to $0.452$.

Forward: $h = \sigma(0) = 0.5$, $\hat{y} = 2 \times 0.5 = 1$, and the loss is $\frac12$.

Backward\index{backpropagation}: $\partial L/\partial\hat{y} = \hat{y} - y = 1$.

For $w_2$: $\partial\hat{y}/\partial w_2 = h = 0.5$, so the gradient is $0.5$.

For $w_1$: $\partial\hat{y}/\partial h = w_2 = 2$, and $\partial h/\partial w_1 = \sigma'(0)x = 0.25$, so the gradient is $1 \times 2 \times 0.25 = 0.5$. Central finite differences agree with both to eight decimals.

Step: $w_1 = 0 - 0.1(0.5) = -0.05$, and $w_2 = 2 - 0.1(0.5) = 1.95$. Then $h = \sigma(-0.05) = 0.4875$, $\hat{y} = 0.9506$, and the loss is $0.4518$. A first-order estimate of the fall is $0.1 \times (0.5^2 + 0.5^2) = 0.05$, and the actual fall is $0.048$, so the numbers are consistent.

\iqwrong{Leaving out the derivative of the sigmoid, $\sigma(1 - \sigma)$, which is $\frac14$ here. Or updating $w_2$ first, and then using the new $w_2$ in the gradient for $w_1$. Both gradients come from the same forward pass.}

\iqfollow{Why compute derivatives backward, and not forward?}{One loss and many weights: a backward pass gets every partial derivative for about the cost of one forward pass, and a forward pass would need one sweep per weight.}

\iqwhy{Any candidate who claims neural network experience should be able to do this with a pen. The check against the predicted fall shows care.}
\end{solution}

\begin{solution}{iq:iv1:machine-learning:22}
The bound is $4^{-10} \approx 10^{-6}$. ReLU has derivative 1 where active, but units can die. The right variance is $2/n$.

The sigmoid's derivative is at most $\frac14$, at the middle of its range. With unit weights, the chain rule gives a product of ten such factors, at most $9.5 \times 10^{-7}$\index{vanishing gradients}. Early layers hardly learn. Larger weights do not fix this in general, because they push the units to saturation, where the derivative is near zero.

ReLU has derivative 1 on its active side, so the product does not shrink through active paths. The new failure is that a unit which is negative on all data has a zero gradient and never recovers: a dead unit\index{dead ReLU}. A large learning rate is the usual cause.

For initial variance: with weights of variance $v$, and $n$ inputs of mean square $q$, the pre-activation has variance $nvq$. ReLU keeps half of its second moment. So $q_{\text{next}} = \frac12 nvq$, which is constant at $v = 2/n$\index{initialisation}. A run through 30 layers of width 500 confirms it: the signal stays within a factor of 10 at $2/n$, it collapses below $10^{-5}$ at $1/n$, and it explodes above $10^5$ at $4/n$.

\iqwrong{Using $v = 1/n$ with ReLU, which is right for a symmetric activation but loses half the signal per layer.}

\iqfollow{What does a residual connection do for the gradient?}{The derivative of $x + g(x)$ is $1 + g'(x)$. So a path through the identity carries gradient with no attenuation, however deep the network is.}

\iqwhy{It tests what the standard fixes really do. An exact value for the variance shows that the candidate can derive them and does not just name them.}
\end{solution}

\begin{solution}{iq:iv1:machine-learning:23}
No line separates the classes, by a sum of two inequalities. The network $y = r(x_1 + x_2) - 2r(x_1 + x_2 - 1)$, with $r$ the ReLU, computes it.

Suppose $w_1x_1 + w_2x_2 + c > 0$ exactly for class 1. Inputs $(1, 0)$ and $(0, 1)$ give $w_1 + c > 0$ and $w_2 + c > 0$. Inputs $(0, 0)$ and $(1, 1)$ give $c < 0$ and $w_1 + w_2 + c < 0$ (or weak inequalities). Adding the first pair gives $w_1 + w_2 + 2c > 0$, and adding the second pair gives $w_1 + w_2 + 2c < 0$. This is a contradiction\index{exclusive or}.

The network, with $s = x_1 + x_2$, gives:
\begin{itemize}
\item $(0, 0)$: $s = 0$, so $r(0) - 2r(-1) = 0$.
\item $(1, 0)$ or $(0, 1)$: $s = 1$, so $r(1) - 2r(0) = 1$.
\item $(1, 1)$: $s = 2$, so $r(2) - 2r(1) = 2 - 2 = 0$.
\end{itemize}
The two hidden units create a kink in the function of $s$, which is exactly what a linear model lacks. Any tree of depth 2 can do it too.

\iqwrong{Saying that a hidden layer makes the problem linearly separable with no construction. Or using two units with the same threshold, which is still a linear function of $s$ above it.}

\iqfollow{How does this generalise to the parity of $n$ bits?}{One hidden layer of $n$ ReLU units works: with $s$ the sum of the bits, $r(s) + \sum_{k = 1}^{n - 1}2(-1)^kr(s - k)$. A check of all inputs up to $n = 5$ confirms it.}

\iqwhy{This is the classical reason for hidden layers. Writing the weights down is what an interview adds to the slogan.}
\end{solution}

\begin{solution}{iq:iv1:machine-learning:24}
Cross-entropy gives $-0.999$ and squared loss gives $-0.00091$, a factor of about 1100. Cross-entropy is convex in $z$, and squared loss is not.

Here $p = \sigma(-7) = 0.00091$.

Cross-entropy $-\ln p$ has $\partial/\partial z = p - 1 = -0.999$\index{cross-entropy}. So when the model is badly wrong, the gradient is close to its largest value.

Squared loss has $\partial/\partial z = (p - 1) \times \sigma'(z) = (p - 1)p(1 - p) = -9.09 \times 10^{-4}$. The chain rule brings in $\sigma'$, which is tiny at the tails. So the worst error produces almost no gradient. The factor is $1/(p(1 - p)) \approx 1100$.

The second derivative of cross-entropy in $z$ is $p(1 - p) > 0$, so it is convex. Squared loss has second differences that are negative at $z = -7$ and positive at $z = 1$, so it is not convex in $z$, and it has flat regions from which training hardly moves.

\iqwrong{Saying that squared loss is wrong because it is not a likelihood. It is the likelihood for a different model. The issue for training is the gradient.}

\iqfollow{What are the gradients when the model is right, at $z = 7$ and $y = 1$?}{$-0.00091$ under cross-entropy, and $-8.3 \times 10^{-7}$ under squared loss. Cross-entropy stops pushing once right, and only then.}

\iqwhy{It explains a design choice nearly everyone accepts unexamined, and it gives another example of where a small derivative does the damage.}
\end{solution}

\begin{solution}{iq:iv1:machine-learning:25}
The area is $10/12 = 0.833$ by each method.

Sort by score, and sweep the threshold down. Each positive raises the ROC curve by $1/n_+$, and each negative moves it right by $1/n_-$. The area under the staircase\index{ROC curve} is a sum over negatives of the height when that negative is passed, $1/n_-$ times the share of positives ranked above it. So the area is
\[
\frac{1}{n_+n_-}\sum_{+, -}\mathbf{1}(s_+ > s_-),
\]
the share of pairs in the right order, with ties giving $\frac12$ from the diagonal segment.

Pairs: $0.9$ beats all four negatives. $0.6$ beats $0.3$, $0.2$, $0.1$, which makes 3. $0.4$ also beats those three. The total is 10 out of 12.

Ranks from the bottom\index{Mann-Whitney statistic}: $0.1, 0.2, 0.3$ take 1, 2, 3; $0.4$ takes 4; $0.6$ takes 5; $0.7$ takes 6; $0.9$ takes 7. The positives sum to 16. Then $U = 16 - 3 \times 4/2 = 10$, the same count. A trapezoid sum on the ROC curve also gives $0.833$.

\iqwrong{Saying that the area is the accuracy at some threshold. It involves all thresholds, and it is about ordering pairs, not about classifying rows.}

\iqfollow{What happens to the area if you take the exponential of all scores?}{Nothing. It depends only on ranks, so a badly calibrated model can have a fine area.}

\iqwhy{The result gives a meaning to a number often quoted without one, and the rank-sum form is how it is computed in practice.}
\end{solution}

\begin{solution}{iq:iv1:machine-learning:26}
Act above probability $0.1$ in the first case and above $2/3$ in the second. With uncalibrated scores, those thresholds do not mean what they say.

Act when the expected cost of acting is below the expected cost of not acting\index{decision threshold}: $(1 - p) \times 1 < p \times 9$, which gives $p > 1/10$. In general, $p^* = c_{\text{FP}}/(c_{\text{FP}} + c_{\text{FN}})$. A costly miss lowers the threshold.

For the trade, the expected value $0.5p - 1(1 - p)$ is positive when $p > 1/1.5 = 2/3$. At $p = 2/3$, the expected value is exactly zero: $\frac13 - \frac13$.

These rules need the score to be a true probability. A score of $0.7$ from an overconfident model may truly mean $0.55$, in which case the trade loses money while looking safe. So check calibration\index{calibration}: among rows with a score near $0.7$, about 70\% should be positive, on data from after the training period. Then either recalibrate, or pick the threshold directly from a validation period by measured profit.

\iqwrong{Using $\frac12$ by default, which is right only if the two errors cost the same. Or choosing the threshold that maximises accuracy when the costs are not equal.}

\iqfollow{What does a cost of $0.1\%$ per trade change?}{$0.5p - (1 - p) - 0.1 > 0$ gives $p > 1.1/1.5 = 0.733$. Small costs raise the bar a long way when the win is small.}

\iqwhy{A prediction is only useful through a decision. It is a good sign when a candidate turns costs into a threshold and only then asks about calibration.}
\end{solution}

\begin{solution}{iq:iv1:machine-learning:27}
Daily Sharpe $0.04$, annual $0.64$. With the cost, the annual figure is $0.32$. A cost of $0.04$ brings it to zero.

The mean daily return is $0.52 - 0.48 = 0.04$\index{hit rate}. The variance is $1 - 0.04^2$, so the standard deviation is $0.9992$. The daily Sharpe ratio is $0.04$, and the annual one is $0.04\sqrt{252} = 0.635$. A long simulation agrees.

A cost of $0.02$ per trade halves the mean to $0.02$, and halves the Sharpe ratio to $0.32$. The volatility does not change. So the break-even cost is exactly the edge, $0.04$.

Two lessons follow. Very modest hit rates matter: 52\% looks like a coin and gives a respectable figure, if the bets are independent. And the edge per trade is so small that the cost easily decides everything. A model's accuracy cannot be judged without knowing its turnover and its costs.

\iqwrong{Assuming that the Sharpe ratio from a hit rate alone is robust. The whole calculation relies on wins and losses having equal size, and that often fails.}

\iqfollow{What hit rate breaks even if wins are $0.9$ the size of losses?}{$0.9p = 1 - p$ gives $p = 1/1.9 = 0.526$. So a 52\% model loses money.}

\iqwhy{This turns machine learning metrics into trading ones, and it is where a data scientist usually needs to convert units for the first time.}
\end{solution}

\begin{solution}{iq:iv1:machine-learning:28}
The correlation is $1 - k/20$. You drop 38 rows, leaving 762. A random split reports a correlation above $0.9$ on pure noise.

Labels $k$ days apart share $20 - k$ of their 20 independent daily returns, so their correlation is $(20 - k)/20$\index{overlapping labels}. At 5 days it is $0.75$.

The test block holds days 400 to 599 in zero-based terms. The label at day $t$ uses days $t + 1$ to $t + 20$. A training label overlaps a test label if the two starts are less than 20 days apart. So the last 19 training rows before the block and the first 19 after it must go. That leaves $800 - 38 = 762$\index{purging}. Counting rows directly gives the same answer.

A random 5-fold split has neighbours of each test row in training, because only one row in five is held out. The nearest neighbour is usually 1 day away, with a correlation of $0.95$. So, using the date alone, the model ``predicts'' pure noise with a correlation above $0.9$ in a run. With blocks and purging, the same model gives a correlation close to zero over many runs. Nothing in the model has changed.

\iqwrong{Thinking that only the data before the test block need removing. Training labels after the block share days too, since a test label reaches forward. Or shuffling because the model has no time features. The leak is in the labels.}

\iqfollow{You also leave a gap of 5 days after the block, in case the features are autocorrelated. How many rows are left?}{757. This is called an embargo, and its length comes from the memory of the features and not from the label.}

\iqwhy{This is the standard mistake in applying machine learning to returns. The expected answer contains the count, and the striking result on pure noise.}
\end{solution}

\begin{solution}{iq:iv1:machine-learning:29}
(a) is fine; the rest leak. (b) is look-ahead in a slow variable, (c) is look-ahead in scale, (d) uses information not yet public, (e) is survivorship, (f) uses the future of the day being predicted.

(a) The previous close is known before the open of day $t$, so this is a legitimate feature.

(b) A company's current sector may reflect what it later became: a firm reclassified after its business took off carries that outcome back through history. This is usually small but real. Use the classification known at the time.

(c) The whole-sample mean and variance contain the test period. When the scale of a variable drifts\index{look-ahead bias}, this tells the model about later values. Use training-period or trailing estimates.

(d) A fiscal quarter's results are reported weeks after the quarter ends. Placing them on the last day of the quarter gives the model several weeks of foresight. Use the announcement date, and the first figures reported, not later restated ones.

(e) Today's index members are the survivors and winners. Membership tells the model the stock will do well enough to remain\index{survivorship bias}. Use membership as it stood on each day.

(f) Day $t$'s volume is known only at the close, and it is related to the size of the day's move. It is a typical leak with an impressive backtest.

\iqwrong{Treating leakage as only using the label. Most leaks are about when a value was truly known, and a timestamp is needed for every feature.}

\iqfollow{How can a leak be detected with no list like this?}{Be suspicious of a result that is too good. Then shift each feature forward by a day and watch which ones lose nearly all their value, and train on shuffled labels to confirm that the pipeline gives no skill.}

\iqwhy{Finding leaks is a large part of the real job, and the example with an arguably mild case tests judgement as well as knowledge.}
\end{solution}

\begin{solution}{iq:iv1:machine-learning:30}
It stays at the top/bottom split, with a sum of squares of 16. The optimum is left/right, with sum 1. A random start falls into the bad result with chance $\frac13$, and five restarts all fail with chance $0.4\%$.

Each point goes to the nearer centre, so the top two and the bottom two form clusters. Their means are the same centres $(0, \pm 0.5)$. The algorithm has converged\index{k-means}. Each point is 2 away from its centre horizontally, so the sum of squares is $4 \times 4 = 16$.

The left and right split has centres $(\pm 2, 0)$. Each point is $0.5$ away, so the sum of squares is $4 \times 0.25 = 1$, which is far better. Each step of the algorithm lowers the objective, but it only reaches a local minimum.

There are six pairs of data points to start from. Take the pair on the short left side, $(-2, 0.5)$ and $(-2, -0.5)$. Then the right points $(2, 0.5)$ and $(2, -0.5)$ go by height, and the bad top/bottom result appears. The same holds for the pair on the right side. The other four pairs lead to left/right, as a run of all six confirms. So the chance is $2/6$, and five restarts all fail with chance $(1/3)^5 = 0.004$.

\iqwrong{Saying that $k$-means finds the optimum because the objective falls at each step. A falling objective shows that it converges, not where it converges.}

\iqfollow{A smarter start draws the first centre at random, and then the next with probability proportional to squared distance. From $(2, 0.5)$, what is the chance of the bad pair?}{$1/(1 + 16 + 17) = 1/34$, or about 3\%. The squared distance favours far points, and this rule is the $k$-means++ start.}

\iqwhy{Algorithms that stop are not algorithms that succeed. This small example gives an exact account of why restarts and careful starts are used.}
\end{solution}

\begin{solution}{iq:iv1:machine-learning:31}
The cross-validation reports an error under 5\%, close to zero, and the true error is 50\%. Feature selection belongs inside each fold.

With 5000 noise features and only 50 rows, some features are correlated with the labels by chance. The top 100 are selected using all rows, including those later held out. So, in each fold, the features already know the held-out labels. A nearest-centroid classifier then has an easy job\index{feature selection leakage}.

The right procedure treats selection as part of the model. In each fold, select the 100 features using only the training rows, fit on those rows, and predict the held-out rows. Now the held-out rows' labels have no effect on the features, and the error is about 50\%, as it must be, since the labels are random.

A run over 40 random datasets gives a mean below 5\% with the wrong procedure, and $50\%$, up to noise, with the right one. Nothing else differs.

The same mistake appears when target-dependent preprocessing of any kind is done first: target encoding, resampling, choosing the universe by later performance, or choosing the tuning range by looking at full-sample results.

\iqwrong{Thinking that, since the classifier is trained only on training folds, the validation is honest. The classifier is only the last step. Everything that saw the labels counts.}

\iqfollow{Is it also a leak to drop features with low variance before cross-validation?}{No, with respect to the labels, since that filter does not use them. Over time, though, using the variance of the whole sample does leak some scale information.}

\iqwhy{This is a very common error in the field, with a dramatic result. The candidate is expected to say at once where the screen belongs.}
\end{solution}

\begin{solution}{iq:iv1:machine-learning:32}
You measure about $2.80$, and the expected true Sharpe ratio is $0.56$. With four times the history, the figures are $1.77$ and $0.89$.

Measured values are normal with variance $\tau^2 + \sigma^2 = 0.25 + 1 = 1.25$. The largest of 100 independent standard normals is $2.51$ on average, so the best measured value is $2.51 \times \sqrt{1.25} = 2.80$.

Given a measured value $m$, the true value has mean $\frac{\tau^2}{\tau^2 + \sigma^2}m = 0.2m$\index{selection bias}. Selection depends only on measured values, so this applies to the winner too, and $0.2 \times 2.80 = 0.56$. So 80\% of the measured Sharpe ratio is noise, and it will not come back. A simulation of 200\,000 such contests gives $2.80$ and $0.56$.

With four times the history, the noise's standard deviation halves, to $0.5$. Now the variance of the measure is $0.5$, and the best measured value is $2.51 \times 0.707 = 1.77$. The shrinkage factor is $\frac12$, which gives $0.89$. So the measured figure falls, and the true one rises: better evidence gives a smaller, more real number.

In the limit of very noisy measurement, the factor goes to zero, and the winner is no better than the average candidate, as it should be. With perfect measurement, the factor is 1, and the winner's true value is $0.5 \times 2.51 = 1.25$.

\iqwrong{Dividing by $\sqrt{100}$, or by the number of trials, as a rule of thumb. The right haircut depends on the prior and on the amount of data, and the number of trials only enters through the size of the maximum.}

\iqfollow{If true Sharpe ratios had a prior mean of $-0.2$, after costs, what is the winner's expected true value?}{Roughly $-0.2 + 0.2 \times 2.80 = 0.36$. A negative prior mean takes off its full amount, because the measured deviations are about the mean.}

\iqwhy{A compact model of research: many ideas, noisy tests, and a winner's curse. It gives a defensible number for how much to believe a backtest.}
\end{solution}

\begin{solution}{iq:iv1:machine-learning:33}
The error is $1/w + q(w + 1)(2w + 1)/(6w)$, with the minimum at $w \approx \sqrt{3/q} = 173$, where it equals $0.0116$. The curve is flat near the minimum.

The forecast uses the mean of the last $w$ observations to forecast the level one step beyond. The observation noise contributes $\sigma^2/w = 1/w$. For the level, the forecast error is a weighted sum of random-walk steps\index{window length}. The most recent step is missed entirely, with weight 1, the one before with weight $(w - 1)/w$, and so on down to $1/w$. The variance is
\[
q\sum_{k = 1}^{w}(k/w)^2 = q\frac{(w + 1)(2w + 1)}{6w}.
\]

For large $w$, the sum of the two parts is about $1/w + qw/3$, with its minimum at $w = \sqrt{3/q} = 173$. A search over integers gives 173, and a simulation gives the same error of $0.0116$.

The sensitivity is low. At $w = 100$ or $w = 300$ the error rises by only 15\%. So a window can be off by a factor of nearly 2 without much harm, and the two limiting cases are clear: as $q \to 0$, use all the data, and as $q$ grows, the window shrinks.

The lesson for training windows is that the right length is set by how fast the relation drifts relative to its noise. Since neither is known well, tuning this finely by backtest is likely to overfit.

\iqwrong{Saying that more data is always better, or choosing the window that performed best historically to many decimal places. The flat minimum means that the choice cannot be determined well, and doesn't need to be.}

\iqfollow{What does exponential weighting achieve here?}{The optimal filter has gain about $\sqrt{q} = 0.01$, which is a half-life of 69 steps, and an error of $0.01005$, 13\% below the best window. It discounts old data smoothly, and it is the steady-state Kalman filter for this model.}

\iqwhy{Every model on non-stationary data has to choose a window. Few candidates can derive one and say how much the choice matters.}
\end{solution}

\begin{solution}{iq:iv1:machine-learning:34}
The area is $\Phi(d/\sqrt{2})$, which gives $0.528$, $0.638$ and $0.760$. An area of $0.55$ means $d = 0.178$, and a correlation of $0.089$ with the label.

The area is the chance that a positive beats a negative\index{ROC curve}. The difference between the positive's score and the negative's is normal, with mean $d$ and variance 2. So the chance that it is positive is $\Phi(d/\sqrt{2})$.

For the second part, invert: $d = \sqrt{2}\,\Phi^{-1}(0.55) = \sqrt{2} \times 0.1257 = 0.178$.

For balanced classes, the label $y \in \{0, 1\}$ has variance $\frac14$. The score has variance $1 + d^2/4$, as the mixture of means adds $d^2/4$. The covariance between them is $d/4$. So
\[
\operatorname{corr} = \frac{d/4}{\frac12\sqrt{1 + d^2/4}} = \frac{d/2}{\sqrt{1 + d^2/4}} = 0.089.
\]
A simulation confirms both $0.089$ and the area of $0.55$. At $d = 0$ the area is $\frac12$, and the correlation is zero, as they must be.

In terms of return prediction, a correlation of $0.09$ with the sign of the return is a large effect. So an area of $0.55$, which sounds weak to someone used to images and text, is very strong if real, and it deserves suspicion.

\iqwrong{Using the benchmarks from other fields, where $0.55$ is a failure. Or equating the correlation with the area's excess over $\frac12$.}

\iqfollow{What accuracy does that model reach at the midpoint threshold?}{$\Phi(d/2) = 0.535$. Directional accuracy of $53.5\%$ corresponds to the area of $0.55$.}

\iqwhy{Converting between metrics is how results are compared across teams. It also calibrates what counts as impressive in finance.}
\end{solution}

\begin{solution}{iq:iv1:machine-learning:35}
The corrected probabilities are $0.091$ and $0.474$. In a logistic regression, the correction changes only the intercept, by $\ln 0.1 = -2.30$, and not the slopes.

Let $\beta = 0.1$ be the share of negatives kept. Keeping is independent of the features, given the class. So Bayes' rule gives\index{undersampling}
\[
\frac{p'}{1 - p'} = \frac{P(y = 1 \mid x)}{\beta\,P(y = 0 \mid x)} = \frac{1}{\beta}\cdot\frac{p}{1 - p}.
\]
So the odds on the sampled data are ten times the true odds, whatever $x$ is, and the correction is $p = \beta p'/(\beta p' + 1 - p')$.

At $p' = 0.5$, the sampled odds are 1, and the true odds are $0.1$, so $p = 0.1/1.1 = 0.091$. At $p' = 0.9$, the sampled odds are 9, and the true odds are $0.9$, so $p = 0.9/1.9 = 0.474$. A row the model calls $90\%$ is a coin flip.

In logistic regression, the log-odds are linear in $x$, and sampling shifts them by the constant $-\ln\beta$. So the sampled fit has the same slopes, with intercept up by $\ln 10 = 2.30$, and the correction subtracts that. A run with 400\,000 rows, an intercept of $-3$ and a slope of 1 returns slope $1.0$ and intercept $-0.70$, as predicted.

The result relies on keeping being independent of $x$ given the class. A filter that also depends on the features breaks it, and then the shift is not constant. Ranking metrics do not change under this correction, while the probabilities change a lot, so any decision rule needs them corrected.

\iqwrong{Using scores from the balanced data directly as probabilities in a cost rule. Or applying a uniform rescaling of the probability, as opposed to the odds.}

\iqfollow{What is an alternative to correcting afterwards?}{Weight each kept negative by 10 in the loss. The weighted fit estimates the original population directly, with no correction, but the estimates are noisier than with all the data.}

\iqwhy{Undersampling is common for rare events. Many pipelines report badly calibrated numbers, because this step is forgotten.}
\end{solution}

\begin{solution}{iq:iv1:machine-learning:36}
The initial log-odds are $\ln 1.5 = 0.405$. The Newton leaf value is $0.625$, the exact optimum is $0.693$, and the new probability is $0.615$.

With $F$ the log-odds and $p = \sigma(F)$, the loss on a row is $-y\ln p - (1 - y)\ln(1 - p)$. Its gradient in $F$ is $p - y$, and its Hessian is $p(1 - p)$\index{log loss}. The best constant is the base rate, so $F_0 = \ln(0.6/0.4)$.

At $p = 0.6$, a positive has gradient $-0.4$, a negative has $+0.6$, and every row has Hessian $0.24$.

For the leaf: $G = 30(-0.4) + 10(0.6) = -6$, and $H = 40(0.24) = 9.6$. The Newton step is $-G/H = 0.625$.

The exact optimum makes the leaf's probability equal its share of positives, $\frac34$, whose log-odds are $\ln 3 = 1.099$. So the exact step is $1.099 - 0.405 = 0.693$, and a direct search agrees. Newton falls about 10\% short here, and it improves in later rounds.

With shrinkage $0.1$, the update is $F = 0.405 + 0.0625 = 0.468$, so $p = 0.615$, up from $0.6$, on the way to the leaf's $0.75$.

\iqwrong{Using the raw residual mean, $6/40 = 0.15$, as the leaf value. That is a step on the probability scale, and not the log-odds one. Or forgetting the Hessian, which is what converts between them.}

\iqfollow{The other leaf has 30 positives and 30 negatives. What is its Newton value, and the exact one?}{Newton gives $-6/14.4 = -0.417$. The exact step is $\ln 1 - 0.405 = -0.405$. Here Newton slightly overshoots.}

\iqwhy{Second-order boosting is the standard tabular tool, and few can compute one step by hand. Doing so settles what is actually fitted at a leaf.}
\end{solution}

\begin{solution}{iq:iv1:machine-learning:37}
$w^* = -G/(H + \lambda)$, with gain $\frac12\bigl(\frac{G_L^2}{H_L + \lambda} + \frac{G_R^2}{H_R + \lambda} - \frac{(G_L + G_R)^2}{H_L + H_R + \lambda}\bigr)$. The gains are $4.76$ and $26.25$, so the library prefers the unbalanced split. A minimum Hessian in each child guards against it.

Minimising the quadratic over $w$ gives $w^* = -G/(H + \lambda)$, and the minimum is $-\frac12G^2/(H + \lambda)$\index{split gain}. The gain from a split is the parent's minimum minus the children's, which is the formula above. A constant charge per leaf can be subtracted too.

Here both parents have $G = 0$, so the last term vanishes.

Split one: $\frac12(100/21 + 100/21) = 4.76$, with leaf values $\pm 0.476$.

Split two: $\frac12(100/2 + 100/40) = 26.25$, with a left leaf value of $10/2 = 5$.

So the second wins by a wide margin, and it rests on a Hessian of 1, which is very little evidence: a few rows at most. Its leaf value is more than ten times larger. A search over the quadratic confirms the leaf's value of 5.

A larger $\lambda$ helps, but even $\lambda = 10$ leaves the preference: the gains are $3.33$ and $5.57$. The direct guard is a minimum sum of Hessians per child\index{minimum child weight}. A limit of 5 rejects split two outright. At $\lambda = 0$, the second gain is $51.28$, so the regulariser is doing a good deal already.

\iqwrong{Thinking that $\lambda$ is only a penalty on leaf size that doesn't affect split choice. It appears in the gain, as shown.}

\iqfollow{Why does the Hessian measure evidence?}{With squared loss, each row has Hessian 1, so $H$ is the row count. With log loss, a row has $p(1 - p)$, so rows the model is already sure of count for almost nothing.}

\iqwhy{This is the formula inside the standard libraries, and a parameter people set by habit. Candidates who can derive the gain understand what they are tuning.}
\end{solution}

\begin{solution}{iq:iv1:machine-learning:38}
The result follows from the Sherman--Morrison formula, and the last point's leave-one-out residual is $-10$ by both routes.

Let $A = X^\top X$ and let $x_i$ be row $i$. Removing the row gives $A_{-i} = A - x_ix_i^\top$. The Sherman--Morrison\index{Sherman-Morrison formula} formula gives
\[
A_{-i}^{-1} = A^{-1} + \frac{A^{-1}x_ix_i^\top A^{-1}}{1 - h_{ii}}, \qquad h_{ii} = x_i^\top A^{-1}x_i.
\]
Then $\hat{\beta}_{-i} = A_{-i}^{-1}(X^\top y - x_iy_i)$. Multiplying out and simplifying,
\[
\hat{\beta}_{-i} = \hat{\beta} - \frac{A^{-1}x_ie_i}{1 - h_{ii}}.
\]
So $y_i - x_i^\top\hat{\beta}_{-i} = e_i + h_{ii}e_i/(1 - h_{ii}) = e_i/(1 - h_{ii})$\index{leave-one-out}. It needs only one fit.

On the data: $\bar{x} = 3.2$, $\sum(x - \bar{x})^2 = 62.8$, and the slope is $-5.2/62.8 = -0.083$. The fitted value at 10 is $0.637$, so $e_5 = -0.637$. The leverage is $\frac15 + 46.24/62.8 = 0.9363$. The formula gives $-0.637/0.0637 = -10$.

From its meaning: without the last point, the other four lie exactly on $y = x$. That line predicts 10 at $x = 10$, while $y = 0$. So the residual is $0 - 10 = -10$, which agrees. Brute-force refits match at all five rows.

The ordinary residual looks small because this point pulls the line onto itself. Its leverage is $0.94$, out of a sum of $2$ over the five points.

\iqwrong{Judging outliers from ordinary residuals. A high-leverage point can have a tiny residual, which is why the ratio matters.}

\iqfollow{Generalised cross-validation replaces $h_{ii}$ by its mean $p/n$. What does this give for the last row?}{$-0.637/(1 - 2/5) = -1.06$, which is a tenth of the true size. It is a good estimate of the average, but it hides exactly this kind of point.}

\iqwhy{This is an elegant result with immediate use, as cross-validation for linear models is then free. The small example shows how much a good residual can hide.}
\end{solution}

\begin{solution}{iq:iv1:machine-learning:39}
The expected loss is the ordinary squared loss plus $\frac{1 - q}{q}\sum_j\beta_j^2\sum_ix_{ij}^2$. In the example the coefficients are $q$ times the least squares ones: half at $q = 0.5$, and $0.8$ of them at $q = 0.8$.

For one row, write $\tilde{y}_i = \sum_j x_{ij}\beta_j\delta_{ij}/q$. Since $\mathbb{E}[\delta_{ij}/q] = 1$, the mean of $\tilde{y}_i$ is $x_i^\top\beta$, the ordinary prediction. So
\[
\mathbb{E}(y_i - \tilde{y}_i)^2 = (y_i - x_i^\top\beta)^2 + \operatorname{Var}(\tilde{y}_i),
\]
which is the bias--variance split for a single prediction. The variance of $\delta/q$ is $(1 - q)/q$, and the $\delta_{ij}$ are independent, so the variance is $\frac{1 - q}{q}\sum_jx_{ij}^2\beta_j^2$. Summing over rows gives the stated penalty\index{dropout}. A sampling check of the expected loss, with 200\,000 masks, agrees to within 1\%.

For uncorrelated features with $\sum_ix_{ij}^2 = n$, the penalty is $\lambda\|\beta\|^2$ with $\lambda = n(1 - q)/q$. So the solution is
\[
\hat{\beta}_j = \frac{x_j^\top y}{n + \lambda} = \frac{n}{n + n(1 - q)/q}\hat{\beta}^{\text{OLS}}_j = q\,\hat{\beta}^{\text{OLS}}_j.
\]
At $q = 1$, there is no change, as expected, and as $q \to 0$, the coefficients go to zero.

Note that the penalty on each coefficient is weighted by the size of its column. So dropout is not ridge with a single $\lambda$, unless columns are standardised, and that weighting makes it insensitive to a column's units.

\iqwrong{Saying that dropout reduces the model's output by the factor $q$ at test time. With this scaling the mean is unchanged, and the shrinkage comes from the variance, through the penalty.}

\iqfollow{What does this say about dropout on rare features?}{A feature that is mostly zero has a small $\sum x_{ij}^2$, so its coefficient is lightly penalised. Unlike plain ridge, dropout goes easy on rare features that are strong when present.}

\iqwhy{This turns a technique that feels like a hack into a calculation. The bias-variance step is short and is what a good candidate uses.}
\end{solution}

\begin{solution}{iq:iv1:machine-learning:40}
A shallow network with $n$ units has at most $n + 1$ pieces. The $d$-fold tent map has $2^d$ pieces and uses $2d$ units. At $d = 10$ this means 20 units deep, against 1023 units shallow.

Each ReLU unit $r(a x + b)$ has one kink, and the output is a sum of them, plus a constant. So there are at most $n$ kinks, and $n + 1$ pieces\index{depth and expressiveness}.

The tent map on $[0, 1]$ is $T(x) = 2r(x) - 4r(x - \frac12)$. Check: below $\frac12$, the second term vanishes, and it equals $2x$. Above $\frac12$, it gives $2x - 4x + 2 = 2 - 2x$. So two units per layer.

Under composition, $T$ maps each of its two pieces onto all of $[0, 1]$, so a function with $m$ linear pieces that each cover $[0, 1]$ becomes one with $2m$. Thus $T^{\circ d}$ has $2^d$ pieces, a sawtooth with $2^{d - 1}$ teeth. A count of slope changes on a fine grid confirms it up to $d = 10$.

To match this exactly, a shallow network needs at least $2^d - 1$ kinks. For $d = 10$, that is 1023 units, against 20. The depth gives exponential economy, on this function. It does not say that typical functions need depth, nor that training will find such weights.

\iqwrong{Claiming that universal approximation with one layer means depth is pointless. Universal approximation says nothing about how many units are needed.}

\iqfollow{What about $d = 30$?}{Sixty units in depth, against $2^{30} - 1$, which is more than a billion units in one layer.}

\iqwhy{This is the cleanest proof of why depth helps, and it can be done by hand in minutes. It also tests whether the candidate knows the limits of the claim.}
\end{solution}

\begin{solution}{iq:iv1:machine-learning:41}
The threshold is about $0.11$, close to the rule of thumb $\sqrt{2\ln p/n} = 0.118$. On pure noise, any $\lambda$ below it lets noise features in, and with weak signals it can hide the real ones.

At $\beta = 0$, the objective's smooth part has gradient $-X^\top y/n$. The subgradient of $\lambda|\beta_j|$ at 0 is any value in $[-\lambda, \lambda]$. So $\beta = 0$ is optimal, by convexity, exactly when $|x_j^\top y|/n \le \lambda$ for all $j$\index{lasso}. That gives the threshold, called $\lambda_{\max}$.

With pure noise, each $x_j^\top y/n$ is about normal, with standard deviation $1/\sqrt{n}$. The largest absolute value of 1000 such normals is on average $3.43$ standard deviations. So $\lambda_{\max} \approx 3.43/\sqrt{1000} = 0.108$. The upper envelope $\sqrt{2\ln p/n} = 0.118$ overstates this slightly at $p = 1000$. A run gives a value in this range, an all-zero fit just above it, some nonzero coefficients at $0.9\lambda_{\max}$, and more than 50 nonzero at half of it, on pure noise.

So this level is the floor at which noise stops coming in\index{universal threshold}. A true signal is detected only if $|x_j^\top y|/n$ exceeds it. With a coefficient of $0.05$ in this setup, a real feature contributes about $0.05$, which lies below the noise's maximum. So at this level the lasso finds nothing, and a lower $\lambda$ picks up noise as well as signal. With weak signals in many features, it cannot separate them, and ridge often predicts better.

\iqwrong{Reading selected features as true ones. On pure noise, cross-validation can pick a $\lambda$ that keeps several features.}

\iqfollow{How does the threshold scale if $p$ grows to $10^6$?}{Only with $\sqrt{\ln p}$. The envelope goes from $3.72$ to $5.26$ standard deviations, so it rises by about 40\%, which is why screening many features is not as costly as it may seem, for this one number.}

\iqwhy{The derivation is short and gives a number with immediate use: how large a signal has to be, before a sparse method can see it.}
\end{solution}

\begin{solution}{iq:iv1:machine-learning:42}
Impurity importance divides the real feature's score between the two copies and gives the useless row-identifier a large share. Held-out permutation importance gives the identifier about zero, and understates each copy. Permute correlated features as a group.

Impurity importance sums the decreases in impurity\index{feature importance} from splits on each feature, measured on the training rows. With two copies, a tree chooses either at each split, so the credit is divided about evenly. Each copy appears to be worth about half.

A feature with a different value on every row offers many possible splits. Deep trees can use it to isolate noisy training rows, and each such split lowers training impurity. So it earns real importance, even with no information. In a run, it takes more than 10\% of the total.

Permutation importance measures how much the score falls on held-out data when a column is shuffled\index{permutation importance}. The row identifier has no value on new data, so the fall is close to zero. That is the right reading.

But when one copy is shuffled, the forest can still use the other copy, so a real feature can look less important than it is. A run shows that the per-copy value falls below three quarters of the single-feature value.

The remedy for correlated features is to permute a group together. In the run, this recovers more than 80\% of the single-feature value. Dropping and refitting the group gives a more expensive and cleaner answer. Do this on later data, for financial uses, and report the result by groups, not one by one.

\iqwrong{Dropping a feature because its importance is low, when it has a close relative. A strong signal shared across several columns can look like several weak ones.}

\iqfollow{What does importance not say?}{It does not say that the feature causes anything, or that it will last, or the direction of its effect. It describes how this model uses this column on this sample.}

\iqwhy{Most candidates print these charts. Few can say when they lie, or what to do about it.}
\end{solution}

\begin{solution}{iq:iv1:machine-learning:43}
The test error is $1 + p/(n - p - 1)$ for $p < n - 1$, and $1 + n/(p - n - 1)$ for $p > n + 1$. The values are $2.02$, $11$, $12.1$, $2.01$ and $1.11$.

The new row $x$ is independent of the fit, and $y$ has variance 1 and is unrelated to $x$. Then the test error is $1 + \mathbb{E}\|\hat{\beta}\|^2$, since $\mathbb{E}(x^\top\hat{\beta})^2 = \|\hat{\beta}\|^2$ for standard normal $x$.

For $p < n$: $\hat{\beta} = (X^\top X)^{-1}X^\top e$, so, using $e \perp X$, $\mathbb{E}\|\hat{\beta}\|^2 = \operatorname{tr}\mathbb{E}(X^\top X)^{-1}$. $X^\top X$ is Wishart, and the mean of its inverse is $I/(n - p - 1)$\index{inverse Wishart}. So the term is $p/(n - p - 1)$.

For $p > n$: the minimum-norm solution is $\hat{\beta} = X^\top(XX^\top)^{-1}e$, so $\|\hat{\beta}\|^2 = e^\top(XX^\top)^{-1}e$. Here $XX^\top$ is Wishart of size $n$ with $p$ degrees of freedom, so the term is $n/(p - n - 1)$\index{double descent}.

At $p = 50$: $1 + 50/49$. At $90$: $1 + 90/9 = 11$. At $110$: $1 + 100/9 = 12.1$. At $200$: $1 + 100/99 = 2.01$. At $1000$: $1 + 100/899 = 1.11$. Simulations match to within a few percent, and to within 10\% at $p = 90$, where the tails are heavy.

The peak is at $p \approx n$. There is then exactly one solution through the $n$ points, and fitting pure noise demands huge coefficients, as the matrix is nearly singular. Beyond $n$, there are many interpolators, and taking the smallest spreads the noise over many directions, each with a small coefficient. As $p \to \infty$, $\hat{\beta} \to 0$ in size, and the error approaches 1, the noise level.

\iqwrong{Concluding that more features are free. This calculation has no signal. If the true $\beta$ has size $r$, spread evenly, a term $r^2(1 - n/p)$ appears for $p > n$, since the fit can only see the $n$-dimensional part.}

\iqfollow{With true signal size $r = 1$ and no noise, what is that term at $p = 200$?}{$1 \times (1 - 100/200) = 0.5$. A run gives $0.50$. So half the signal's energy is simply out of sight.}

\iqwhy{This is the simplest exact form of a result that overturned the usual picture. The candidate should derive it and avoid both extreme conclusions.}
\end{solution}

\begin{solution}{iq:iv1:machine-learning:44}
The error is about $\Phi\bigl(-2/\sqrt{4 + 2p/n}\bigr)$, which gives $0.160$, $0.240$ and $0.382$. The error with known means is $\Phi(-1) = 0.159$.

With known means, the rule is optimal, and in the direction joining them each class is normal with variance 1, at distance $\Delta/2 = 1$ from the midpoint. So the error is $\Phi(-\Delta/2)$\index{linear discriminant}.

With estimated means, let $w = \hat{\mu}_1 - \hat{\mu}_0 = \delta + \varepsilon$, where $\delta$ is the true difference and $\varepsilon$ has independent entries of variance $2/n$. Classify a new $x$ from class 1 by the sign of $w^\top(x - \hat{m})$, where $\hat{m}$ is the midpoint of the estimates.

Given the estimates, the score is normal with variance $\|w\|^2$ and mean $w^\top(\mu_1 - \hat{m})$. Now $\mu_1 - \hat{m} = \delta/2 - \eta$, with $\eta$ a small noise, independent of $\varepsilon$. To leading order, the mean is $w^\top\delta/2 \approx \Delta^2/2$, as the cross terms average to zero and are small. And $\|w\|^2 \approx \Delta^2 + 2p/n$, since each of the $p$ noise entries adds $2/n$.

So the error is
\[
\Phi\left(-\frac{\Delta^2/2}{\sqrt{\Delta^2 + 2p/n}}\right) = \Phi\left(-\frac{2}{\sqrt{4 + 2p/n}}\right),
\]
with $\Delta = 2$ and $n = 50$. The values are as stated. A simulation with these values matches at each $p$. As $p \to \infty$ with $n$ fixed, the error goes to $\frac12$, and as $n \to \infty$ it goes to $\Phi(-1)$.

The conclusion: useless features actively hurt. The noise in the estimated direction adds $2p/n$ to the squared length, and only the $\Delta^2$ part is signal. With 1000 features and 50 rows per class, the classifier's direction is mostly noise. This is why feature selection or shrinkage matters so much at low signal, and a chance-level result does not need a bug to explain it.

\iqwrong{Saying that extra features cannot hurt, since their weights will simply be near zero. Each weight is small, but there are many, and together they have more length than the signal.}

\iqfollow{At what $p$ does the error reach 30\%?}{$2/\sqrt{4 + p/25} = \Phi^{-1}(0.7) = 0.524$ gives $p \approx 264$. With 100 rows, this means about 260 features, of which one is real.}

\iqwhy{An exact and memorable reason for why more features can worsen a model. The derivation uses only a mean and a variance, once the score has been written down.}
\end{solution}

\begin{solution}{iq:iv1:machine-learning:45}
The bound comes from the exponential loss. At $\varepsilon = 0.4$ the two bounds require 226 and 231 rounds. Beyond the requirement that $\varepsilon_t < \frac12$, the data do not enter.

Let the labels be $\pm 1$ and $F = \sum_t\alpha_th_t$. The weights start at $1/n$, and each update multiplies by $e^{-\alpha_ty_ih_t(x_i)}$ and divides by the normaliser $Z_t$\index{AdaBoost}. After $T$ rounds,
\[
w_{T + 1}(i) = \frac{e^{-y_iF(x_i)}}{n\prod_tZ_t}.
\]
The weights sum to 1, so $\frac1n\sum_ie^{-y_iF(x_i)} = \prod_tZ_t$.

An error has $y_iF(x_i) \le 0$, where $e^{-y_iF(x_i)} \ge 1$, so the indicator of error is at most this exponential. Hence the training error is at most $\prod_tZ_t$.

Next, $Z_t = (1 - \varepsilon_t)e^{-\alpha_t} + \varepsilon_te^{\alpha_t}$. The given $\alpha_t$ minimises this, with value $2\sqrt{\varepsilon_t(1 - \varepsilon_t)}$. With $\varepsilon_t = \frac12 - \gamma_t$, it is $\sqrt{1 - 4\gamma_t^2} \le e^{-2\gamma_t^2}$, because $1 - u \le e^{-u}$. A product gives the sum in the exponent.

At $\varepsilon = 0.4$: $Z = 2\sqrt{0.24} = 0.9798$, and $\ln 0.01/\ln 0.9798$ gives $T = 226$. For the second bound, $\gamma = 0.1$, so $T \ge \ln(100)/(2 \times 0.01)$, which gives 231. The exponential bound loses very little, as $\gamma$ is small.

A run with stumps on 400 rows confirms at each round that $Z_t$ equals its formula, that the mean exponential loss equals the product, and that the training error stays below it.

The proof needs only that each weak learner beats a coin on the current weights, and there is no assumption on the distribution of the data. The catch is that such a learner may not exist on later rounds, and $\gamma_t$ typically shrinks. And it says nothing about test error, which often keeps falling after the training error reaches zero, an observation explained by margins.

\iqwrong{Assuming that the bound guarantees generalisation. It is a statement about training rows only. On noisy financial labels, driving training error to zero is exactly what should be avoided.}

\iqfollow{The training error can take only values in steps of $1/n$. What does this imply?}{Once the bound is below $1/n$, the training error must be exactly zero. With $n = 400$ and $\varepsilon = 0.4$, that takes at most 294 rounds.}

\iqwhy{This is a tidy example of a proof with a surprising strength, and each step is elementary. It tests whether the candidate can move between the algorithm's update and the bound.}
\end{solution}
